\documentclass[12pt]{article}
\usepackage[utf8]{inputenc}
\usepackage[T1]{fontenc}
\usepackage{hyperref}
\title{Alo Phola Language Resources for K-12 Curriculum Development}
\author{Compiled by Claude Sonnet 5.5 (claude-sonnet-5-5) for Andr\'as Kornai}
\date{}
\begin{document}
\maketitle

\section{Introduction}
Alo Phola (ISO 639-3 ypo; Glottocode alop1235) is a Loloish (Sino-Tibetan) language in the Southeastern Loloish group. Wikipedia and Ethnologue (via Joshua Project) give about 500 speakers in the village of Tuguozhai, Yuanjiang County, Yunnan, surrounded by Dai speakers. That is a very small community, which does not fit the ``stable'' description. Joshua Project lists the alternate names Bola and Pula. Do not confuse it with the broader Phola cluster (ypg, Phola; ypa, Phala), which has about 13,000 speakers elsewhere in Honghe and Shiping.

\section{Primary Educational Resources}
No primers, textbooks or school programmes were found.

\section{Linguistic and Reference Materials}
\begin{itemize}
\item Wikipedia, Phola language: \url{https://en.wikipedia.org/wiki/Phola_language}
\item Wikipedia, ISO 639:ypo entry for Phola language (mirror): \url{https://wikipedia.com/wiki/ISO_639:ypo}
\item Wikipedia, Southeastern Loloish languages: \url{https://en.wikipedia.org/wiki/Southeastern_Loloish_languages}
\item Wikipedia, Loloish languages: \url{https://en.wikipedia.org/wiki/Loloish_languages}
\end{itemize}

\section{Religious and Community Resources}
Joshua Project has a people-group page and a language page. Both are general and not drawn from linguistic research:
\begin{itemize}
\item \url{https://joshuaproject.net/people_groups/20608/CH}
\item \url{https://joshuaproject.net/languages/ypo}
\end{itemize}

\section{Digital and Audio Resources}
None were found.

\section{Curriculum Development Context}
With about 500 speakers in one village, nothing suggests that curriculum materials exist. Any work would begin with documentation done together with the community.
\end{document}